\subsection{验证 Rim--Schlessinger 判别法}\label{src:C.4.3-subsection}
我们通过验证 $(\mathrm{RS}_1)$--$(\mathrm{RS}_3)$, 应用 Rim--Schlessinger 判别法 (\bookref{C.4.6}), 为三个主要模问题构造极小万有形式形变. 实际上, 我们将验证以下\textit{强齐性}条件:
\begin{enumerate}[(1)]
\item[$(\mathrm{RS}_4^*)$] 对 (未必为 Artin 的) 环的每个映射 $B_0\to A_0$ 及核的平方为零的满射 $A\twoheadrightarrow A_0$, $\mathcal X(B_0\times_{A_0}A)\to\mathcal X(B_0)\times_{\mathcal X(A_0)}\mathcal X(A)$ 为等价.
\end{enumerate}
验证 $(\mathrm{RS}_4^*)$ 常与验证较弱的条件 $(\mathrm{RS}_1)$--$(\mathrm{RS}_2)$ 同样容易. 强齐性还将作为第二版 Artin 代数性公理 (\bookref{C.7.4}) 中的一条公理, 因为它有助于验证万有性的开性. 实际上每个代数叠都满足 $(\mathrm{RS}_4^*)$ (见 \cite[Tag 07WN]{SP}); 换言之, Ferrand 推出 $\Spec(B_0\times_{A_0}A)$ 是代数叠范畴中的推出.

验证 $(\mathrm{RS}_4^*)$ 依赖环的纤维积上的模的性质.
\begin{lemma}\label{src:C.4.12}
设 $A\twoheadrightarrow A_0$ 为核的平方为零的环满射, $B_0\to A_0$ 为环映射. 设 $M,M_0,N_0$ 分别为 $A,A_0,B_0$ 上的平坦模, $M\to M_0$ 为 $A$-模映射, $N_0\to M_0$ 为 $B_0$-模映射. 假设 $M\otimes_AA_0\to M_0$ 及 $N_0\otimes_{B_0}A_0\to M_0$ 均为同构. 令 $B:=B_0\times_{A_0}A$, $N=N_0\times_{M_0}M$. 则
\begin{enumerate}[(1)]
\item 映射 $N\otimes_BB_0\to N_0$ 和 $N\otimes_BA\to M$ 均为同构,
\item $N$ 在 $B$ 上平坦, 且
\item 模 $N_0$ 与 $M$ 有限表现, 当且仅当 $N$ 有限表现.
\item 设 $B_0'$ 为 $B_0$-代数, $A'$ 为 $A$-代数, $B_0'\otimes_{B_0}A_0\to A'\otimes_AA_0$ 为同构. 令 $B'=B_0'\times_{A_0'}A$. 则 $B\to B'$ 平坦且有限表现, 当且仅当 $B_0\to B_0'$ 及 $A\to A'$ 均如此.
\end{enumerate}
\end{lemma}
\begin{proof}
我们验证 (2), 其余断言留给读者. 令 $J=\ker(A\to A_0)$. 因为 $M$ 在 $A$ 上平坦, 局部平坦性判别法 (\bookref{A.2.6}) 蕴含 $0\to J\otimes_{A_0}M_0\to M\to M_0\to0$ 正合, 故
\begin{equation}\label{src:C.4.13}
0\to J\otimes_{A_0}M_0\to N\to N_0\to0
\end{equation}
也正合. 为证明 $N$ 平坦, 由局部平坦性判别法 (\bookref{A.2.6}), 只需验证 (a) $N\otimes_BB_0$ 在 $B_0$ 上平坦, 以及 (b) $J\otimes_{B_0}N_0\to N$ 单射. 由 (2), (a) 中的模与 $N_0$ 等同, 后者由假设在 $B_0$ 上平坦. 因为 $J\otimes_{A_0}M_0=J\otimes_{A_0}(N_0\otimes_{B_0}A_0)=J\otimes_{B_0}N_0$, (b) 由 \bookeqref{C.4.13} 推出. 见 \cite[引理 $3.4$]{Sch68}, \cite[命题 $16.4$]{Har10} 及 \cite[Tags 08KG, 0D2G 和 0D2K]{SP}.
\end{proof}
若对每个 $x_0\in\mathcal X(\Bbbk)$, 局部形变预叠 $\mathcal X_{x_0}$ (定义见注记 \bookref{C.4.4}) 具有极小万有形式形变 (相应地, 原可表, 满足 $(\mathrm{RS}_3)$), 则称 $\Sch/\Bbbk$ 上的预叠 $\mathcal X$ \textit{具有极小万有形式形变} (相应地, \textit{局部原可表}, \textit{满足 $(\mathrm{RS}_3)$}).
\begin{proposition}\label{src:C.4.14}
设 $X$ 为域 $\Bbbk$ 上的固有概形.
\begin{enumerate}[(1)]
\item Hilbert 函子 $\Hilb(X):\Sch/\Bbbk\to\Sets$, 其在 $S$ 上的对象为在 $S$ 上平坦且有限表现的闭子概形 $Z\subseteq X_S$, 满足 $(\mathrm{RS}_3)$ 及 $(\mathrm{RS}_4^*)$, 因此局部原可表.
\item $\Sch/\Bbbk$ 上的叠 $\mathcal{F}am$, 其在 $S$ 上的对象为代数空间的固有, 平坦且有限表现态射 $\mathcal Y\to S$,\footnote{若要使 $\mathcal{F}am$ 为叠, 则必须允许 $\mathcal Y$ 为代数空间; 见例 \bookref{3.5.12} 和习题 \bookref{5.5.15}. 另一方面, Rim--Schlessinger 判别法 (\bookref{C.4}) 同样适用于参数化概形的固有, 平坦且有限表现态射 $\mathcal Y\to S$ 的预叠.} 具有极小万有形式形变. 特别地, 所有曲线组成的叠 $\mathcal M_g^{\mathrm{all}}$ 和光滑曲线叠 $\mathcal M_g$ 具有极小万有形式形变.
\item $\Sch/\Bbbk$ 上的叠 $\Coh(X)$, 其在 $S$ 上的对象为 $S$ 上平坦的有限表现拟凝聚 $\mathcal O_{X_S}$-模, 满足 $(\mathrm{RS}_3)$ 和 $(\mathrm{RS}_4^*)$, 因此具有极小万有形式形变. 特别地, 若 $C$ 为光滑连通射影曲线, 则 $\cBun(C)$ 具有极小万有形式形变.
\end{enumerate}
\end{proposition}
\begin{proof}
对于 (1), 命题 \bookref{C.1.3} 将 $\Hilb(X)$ 在 $Z_0\subseteq X$ 处的切空间等同于 $\mathrm H^1(Z,N_{Z_0/X})$. 因为 $X$ 固有, $\mathrm H^1(Z,N_{Z_0/X})$ 有限维, 所以 $(\mathrm{RS}_3)$ 成立. 为检验 $(\mathrm{RS}_4^*)$, 设 $A\twoheadrightarrow A_0$ 为核的平方为零的环满射, $B_0\to A_0$ 为环映射, 假设 $W_0\subseteq X_{B_0}$ 及 $Z\subseteq X_A$ 为基底上平坦的闭子概形, 满足 $Z_0:=W_0\times_{B_0}A_0=Z\times_AA_0\subseteq X_{A_0}$. 则 $\mathcal O_{W_0}\times_{\mathcal O_{Z_0}}\mathcal O_Z$ 为 $\mathcal O_{X_B}$ 的商, 定义闭子概形 $W\subseteq X_B$. 对每个仿射开子集 $\Spec B'\subseteq X_B$, 其限制为 $\Spec B_0'\subseteq X_{B_0}$, $\Spec A'\subseteq X_A$ 及 $\Spec A_0'\subseteq X_{A_0}$, 有 $B'=B_0'\times_{A_0'}A'$. 由引理 \bookref{C.4.12}(4), $B\to B'$ 平坦且有限表现, 因而 $W$ 在 $B$ 上平坦且有限表现.

对于 (2), 由定理 \bookref{C.3.6}, $\mathcal{F}am$ 在 $Y_0$ 处的切空间等同于 $\operatorname{Ext}_{\mathcal O_{Y_0}}^1(Y_0,\mathcal O_{Y_0})$, 因而有限维. (当 $Y_0$ 光滑时, 切空间为 $\mathrm H^1(Y_0,T_{Y_0})$.) 若 $[\mathcal Z_0\to\Spec B_0]\in\mathcal{F}am(B_0)$ 及 $[\mathcal Y\to\Spec A]\in\mathcal{F}am(A)$, 带有同构 $(\mathcal Z_0)_{A_0}\xrightarrow{\sim}\mathcal Y_{A_0}$, 则由定理 \bookref{B.4.1}, Ferrand 推出 $\mathcal Z:=\mathcal Z_0\amalg_{\mathcal Y_{A_0}}\mathcal Y$ 存在. 对 $\mathcal Z$ 的仿射覆盖应用引理 \bookref{C.4.12}(4), 可知 $\mathcal Z\to\Spec B$ 平坦且有限表现. 此外, 因为 $\mathcal Z$ 为推出, $\mathcal Z_0$ 与 $\mathcal Y$ 的相容同构唯一延拓为 $\mathcal Z$ 的同构.

对于 (3), 由命题 \bookref{C.1.18}, $\Coh(X)$ 在凝聚层处的切空间等同于 $\operatorname{Ext}_{\mathcal O_X}^1(E,E)$, 它有限维. 基变换 $X_B$ 为 $X_{A_0}\to X_{B_0}$ 与 $X_{A_0}\hookrightarrow X_A$ 的推出; 对每个仿射开子集 $\Spec B'\subseteq X_B$, 其限制为 $\Spec B_0'\subseteq X_{B_0}$, $\Spec A'\subseteq X_A$ 和 $\Spec A_0'\subseteq X_{A_0}$, 有 $B'=B_0'\times_{A_0'}A'$. 对限制到 $F_0\in\Coh(X)(A)$ 的 $G_0\in\Coh(X)(B_0)$ 和 $F\in\Coh(X)(A)$, 在 $X_B$ 上定义 $G:=G_0\times_FF$. 由引理 \bookref{C.4.12}(3) (对 $B'=B_0'\times_{A_0}A'$ 上的 $G|_{\Spec B'}$ 应用), 它有限表现; 由引理 \bookref{C.4.12}(2) (对 $B=B_0\times_{A_0}A$ 上的 $G|_{\Spec B'}$ 应用), 它在 $B$ 上平坦. 这给出 $(\mathrm{RS}_4^*)$ 的本质满射性, 全忠实性则显然.
\end{proof}
\begin{exercise}\label{src:C.4.15}
若 $X_0$ 为 $\Bbbk$ 上没有无穷小自同构的光滑固有概形, 即 $\mathrm H^0(X_0,T_{X_0})=0$, 证明 $X_0$ 的形变函子原可表.
\end{exercise}
\begin{exercise}\label{src:C.4.16}
设 $X=\Spec A$ 为域 $\Bbbk$ 上具有孤立奇点的仿射概形, 设 $F_X:\mathrm{Art}_{\Bbbk}\to\Sets$ 为如下函子: $F_X(R)$ 是二元组 $(B,\beta)$ 的同构类集合, 其中 $B$ 为平坦 $R$-代数, $\beta:A\to B\otimes_RR/\mathfrak m_R$ 为同构. 证明 $F_X$ 具有极小万有形式形变.
\end{exercise}
